\emph{Conventions.} $T_{i}$ is the translation by one site along the axis $i$; $s_{i} = T_{i} - 1$, $b_{i} = 1 - T_{i}^{-1}$, so that $\Delta_{i} = b_{i}s_{i} = T_{i} - 2 + T_{i}^{-1}$; all of these commute. For the sum over sites the adjoint of $s_{i}$ is $-b_{i}$ and that of $b_{i}$ is $-s_{i}$. Every lattice operator below is a matrix of polynomials in the $T_{i}^{\pm 1}$, and every identity of (2) is an identity between such polynomials; each is also verified on sample fields in the check.

\emph{(1).} The stored proof of Theorem \claimref{synthesis.first-order-gravitational-field-has-two-free-directions-at-null-momentum-and-four-constraints} gives, with $u = h(k)$, $\tau = \tr h$, $Q = Q(k)$,
\[
\mathcal{G}_{k}[h](a) = -\tfrac12\bigl[(a\cdot k)(u - \tau k) - Q\,h(a) + (a\cdot u)\,k\bigr] + \tfrac12\,(k\cdot u - Q\tau)\,a .
\]
Let $h(e_{0}) = 0$ and $k = \omega e_{0} + \kappa$. Then $u = h(\kappa)$ is orthogonal to $e_{0}$, $k\cdot u = \kappa\cdot h(\kappa)$ and $Q = \omega^{2} + Q(\kappa)$. For $a = e_{0}$: $a\cdot k = \omega$, $h(a) = 0$, $a\cdot u = 0$, so $\mathcal{G}_{k}[h](e_{0}) = -\tfrac12\omega\,(h(\kappa) - \tau\kappa) + \tfrac12\omega^{2}\tau\,e_{0} + \tfrac12(\kappa\cdot h(\kappa) - \omega^{2}\tau - Q(\kappa)\tau)\,e_{0}$: its component along $e_{0}$ is $\tfrac12(\kappa\cdot h(\kappa) - Q(\kappa)\tau)$, with no $\omega$, and the rest is $-\omega\,\Pi[h](\kappa)$, $\Pi[h] = \tfrac12(h - \mathrm{id}\,\tau)$. For $a$ orthogonal to $e_{0}$: $a\cdot k = a\cdot\kappa$, and collecting the terms with $\omega^{2}$, namely $\tfrac12\omega^{2}h(a) - \tfrac12\omega^{2}\tau a = \omega^{2}\Pi[h](a)$, the part orthogonal to $e_{0}$ is $\omega^{2}\Pi[h](a) - G_{\kappa}[h](a)$ with $G_{\kappa}$ as stated. A stencil that is a product of two differences annihilates constants and linear functions, and on a plane wave of wave vector $\kappa$ its leading term in the lattice spacing is minus one half of its second moment contracted with $\kappa$ twice; a single difference annihilates constants and its leading term is its first moment contracted with $\kappa$. Writing $h(a) = \sum_{m,n}H_{mn}(a\cdot e_{m})e_{n}$, the three expressions above have components $-\tfrac12 D_{\kappa}H$, $\omega\,C_{\kappa}[\Pi]$ and $\omega^{2}\Pi - G_{\kappa}H$ with $\Pi = \tfrac12(H - \delta\,\mathrm{tr}H)$ and $D_{\kappa}$, $C_{\kappa}$, $G_{\kappa}$ the leading symbols of the lattice stencils; the leading symbol does not tell $s_{i}$ from $b_{i}$, so this fixes the lattice operators to leading order only, and their exact form is fixed by the identities of (2). It is a comparison of polynomials of degree two in $\kappa$ and one in $H$, verified in the check at three momenta on a sample $H$ and term by term from the formulas of the statement (each difference $s_{i}$ or $b_{i}$ contributes $\kappa_{i}$, and the product of two contributes $-\kappa_{i}\kappa_{j}$). For a source $\mathcal{T}(a) = \rho\,(a\cdot e_{0})e_{0}$ the stored equation $\mathcal{G}_{k}[h] = \lambda\mathcal{T}$ then reads: $-\tfrac12 Dh = \lambda\rho$ (so $g$ is $-2\lambda$ times the mass); $C_{i}[\pi] = 0$ with $\pi = \Pi[h'] = \tfrac12(h' - \delta\,\tr h')$, a factor $\omega$ being one derivative in the time of the frame, and then $\tr\pi = -\tr h'$ and $h' = 2p$; and $\pi' = -G[h]$, a factor $\omega^{2}$ being minus two derivatives.

\emph{(2).} $D\delta_{\xi}$: the coefficient of $\xi_{i}$ is $-2b_{i}\sum_{k\neq i}\Delta_{k} + 2\sum_{j\neq i}b_{i}b_{j}s_{j} = 0$. $G[\delta_{\xi}]_{ii} = \Delta_{k}b_{j}\xi_{j} + \Delta_{j}b_{k}\xi_{k} - b_{j}b_{k}(s_{j}\xi_{k} + s_{k}\xi_{j}) = 0$, and in $G[\delta_{\xi}]_{ij}$ the coefficients of $\xi_{j}$, $\xi_{i}$, $\xi_{k}$ are $\tfrac12(-\Delta_{k}s_{i} + s_{i}b_{k}s_{k})$, $\tfrac12(-\Delta_{k}s_{j} + s_{j}b_{k}s_{k})$ and $\tfrac12(s_{i}b_{k}s_{j} + s_{j}b_{k}s_{i} - 2s_{i}s_{j}b_{k})$, all zero. Adjoint: $\langle\pi, \delta_{\xi}\rangle = \sum_{x}\bigl[\sum_{i}\pi_{ii}\,2b_{i}\xi_{i} + 2\sum_{i<j}\pi_{ij}(s_{i}\xi_{j} + s_{j}\xi_{i})\bigr] = -2\sum_{x}\sum_{i}\xi_{i}\bigl(s_{i}\pi_{ii} + \sum_{j\neq i}b_{j}\pi_{ij}\bigr)$. Symmetry of $G$: with the weights of the pairing, the entries are $\tfrac12\Delta_{k}$ between $h_{ii}$ and $h_{jj}$ (self-adjoint), $-b_{j}b_{k}$ from $h_{jk}$ to $G_{ii}$ against $2\cdot(-\tfrac12)s_{j}s_{k}$ from $h_{ii}$ to $G_{jk}$ (adjoint of $b_{j}b_{k}$ is $s_{j}s_{k}$), $-\Delta_{k}$ on $h_{ij}$, and $s_{i}b_{k}$ from $h_{jk}$ to $G_{ij}$ against $s_{k}b_{i}$ from $h_{ij}$ to $G_{jk}$ (adjoint of $s_{i}b_{k}$ is $b_{i}s_{k}$). Then $\langle G[h], \delta_{\xi}\rangle = \langle h, G[\delta_{\xi}]\rangle = 0$ for all $\xi$ gives $C_{i}[G[h]] = 0$. Trace: $\sum_{i}G[h]_{ii} = \tfrac12\sum_{i}\sum_{k\neq i}\Delta_{k}h_{ii} - \sum_{j<k}b_{j}b_{k}h_{jk} = -\tfrac12 Dh$. Next $Dp = 2\sum_{i<j}b_{i}b_{j}\pi_{ij} - \sum_{i}\sum_{k\neq i}\Delta_{k}\pi_{ii} + \Delta\,\tr\pi = 2\sum_{i<j}b_{i}b_{j}\pi_{ij} + \sum_{i}\Delta_{i}\pi_{ii} = \sum_{i}b_{i}C_{i}[\pi]$. The constraints: $C_{0}(s)$ is a function of $h$ and the matter, $C_{i}$ of $\pi$; two $C_{0}$ commute (the masses are commuting numbers of the arrangements), two $C_{i}$ commute, and $[C_{0}(s), C_{i}(l)]$ is, by the adjoint identity, $D$ applied to a displacement, zero. A step changes $C_{i}[\pi]$ by $-\tau C_{i}[G[h]] = 0$ and $Dh$ by $2\tau Dp = 2\tau\sum_{i}b_{i}C_{i}[\pi(n+1)] = 0$ on fields with $C_{i} = 0$, and the matter is at rest. A step is the composition of $\pi \mapsto \pi - \tau G[h]$ and $h \mapsto h + 2\tau p$, each the flow of a quadratic function of one of the two variables with a symmetric form, hence canonical. The decomposition of $G$: with $V_{i} = C_{i}[h] - \tfrac12 s_{i}\tr h$, expanding $-\tfrac12\Delta h + \tfrac12\delta_{V} - \tfrac12\delta\,Dh$ component by component gives the two formulas for $G$ (for $G_{ii}$: $-\tfrac12\Delta h_{ii} + b_{i}(s_{i}h_{ii} + \sum_{j\neq i}b_{j}h_{ij} - \tfrac12 s_{i}\tr h) - \tfrac12 Dh$, in which the coefficient of $h_{ii}$ is $-\tfrac12\Delta + \Delta_{i} - \tfrac12\Delta_{i} + \tfrac12\sum_{k\neq i}\Delta_{k} = 0$, that of $h_{jj}$ is $-\tfrac12\Delta_{i} + \tfrac12(\Delta_{i} + \Delta_{k}) = \tfrac12\Delta_{k}$, that of $h_{ij}$ is $b_{i}b_{j} - b_{i}b_{j} = 0$ and that of $h_{jk}$ is $-b_{j}b_{k}$; likewise for $G_{ij}$). If $V = 0$ and $Dh = 0$ then $G[h] = -\tfrac12\Delta h$ and $\tr G[h] = 0$, and $h(n+1) - 2h(n) + h(n-1) = 2\tau(p(n+1) - p(n)) = -2\tau^{2}\bigl(G[h(n)] - \tfrac12\delta\,\tr G[h(n)]\bigr) = \tau^{2}\Delta h(n)$. The count on $27$ sites is a computation of ranks: $162$ components, the displacements have rank $78$ (the three constant $\xi$ act trivially) and $D$ has rank $26$ (its values sum to zero), leaving $58$; at zero momentum $D$ and the displacements vanish and all six constant fields remain, and the other $52$ are two for each of the $26$ nonzero momenta, since each nonzero momentum contributes at least $6 - 3 - 1 = 2$ and the total is $52$ (a pair of real directions for each momentum, as the bivectors of the stored photon mode are in Proposition \claimref{field-theory.fock-space-maxwell-modes}(2)).

\emph{(3).} For a function $P$ on the sites and a stencil $O$ with reflected stencil $O^{\dagger}$, $\sum_{x}P(x)\,(Of')(x) = \sum_{y}f'(y)\,(O^{\dagger}P)(y)$, and the right side involves $P$ only within one site of the support of $f'$. The stencils of $D$ are products of two differences, so $O^{\dagger}P = 0$ where $P$ is constant or linear; on a box on which the coordinates are single-valued, with the margin, $P = 1$ and $P = x_{k}$ are such functions. Hence the two sums vanish. $\delta_{a} - \delta_{b}$ has $\sum_{x}x_{k}(\delta_{a} - \delta_{b})(x) = a_{k} - b_{k}$, not zero for some $k$. For the pair, in coordinates from the centre along the axis, the wide density is $\delta_{-2} + \delta_{2}$ and the narrow one $\delta_{-1} + \delta_{1}$: monopoles $2$ and $2$, dipoles $0$ and $0$, second moments $8$ and $2$. With $M$ the indicator of the string, $f_{11} = \tfrac12 M$, $f_{22} = f_{33} = -\tfrac12 M$: $Df = -\bigl[(\Delta_{2} + \Delta_{3})\tfrac12 - (\Delta_{1} + \Delta_{3})\tfrac12 - (\Delta_{1} + \Delta_{2})\tfrac12\bigr]M = \Delta_{1}M = \delta_{-2} - \delta_{-1} - \delta_{1} + \delta_{2}$. The last sentence: $W(f')$ and $C_{i}[\pi]$ are both functions of $\pi$. (For the Reading: $p = \tfrac12 f$ is $p[\pi]$ for $\pi_{11} = \tfrac12 M$, all other components zero, since then $\tr\pi = \tfrac12 M$ and $\pi - \tfrac12\delta\,\tr\pi$ has diagonal $(\tfrac14, -\tfrac14, -\tfrac14)M$; and $C_{1}[\pi] = \tfrac12 s_{1}M$ is $+\tfrac12$ on the link from $-2$ to $-1$, $-\tfrac12$ on the link from $1$ to $2$ and zero elsewhere, $C_{2}[\pi] = C_{3}[\pi] = 0$.)

\emph{(3a).} Let $S = G[f']$ vanish outside the box $B$, and take coordinates $x_{m}$ on $B$ enlarged by one site. Let $\xi$ be any function on the links that is given on the enlarged box and is zero elsewhere; $\delta_{\xi}$ at a site of $B$ involves $\xi$ only on links within one site. By (2), $\langle S, \delta_{\xi}\rangle = -2\sum\xi_{i}\,C_{i}[S] = 0$, since $C_{i}[G[f']] = 0$ identically. Three choices. (a) $\xi_{j} = \tfrac12\sum_{i}c_{ij}x_{i}$ with $c$ a constant symmetric matrix: $(\delta_{\xi})_{ii} = 2b_{i}\xi_{i} = c_{ii}$ and $(\delta_{\xi})_{ij} = s_{i}\xi_{j} + s_{j}\xi_{i} = c_{ij}$, the constant field $c$; so $S$ pairs to zero with every constant field, in particular $\sum\tr S = 0$. (b) $\xi_{i} = 2x_{l}x_{i} - \delta_{il}\sum_{m}x_{m}^{2}$, the value on the link from $x$ to $x + e_{i}$ being taken at $x$. For $i \neq l$: $2b_{i}\xi_{i} = 4x_{l}$. For $i = l$: $\xi_{l} = x_{l}^{2} - \sum_{m\neq l}x_{m}^{2}$ and $2b_{l}\xi_{l} = 2(x_{l}^{2} - (x_{l} - 1)^{2}) = 4x_{l} - 2$. For $i, j \neq l$ distinct: $s_{i}\xi_{j} + s_{j}\xi_{i} = 0$. For $i \neq l$: $s_{i}\xi_{l} + s_{l}\xi_{i} = -((x_{i} + 1)^{2} - x_{i}^{2}) + 2x_{i} = -1$. So $\delta_{\xi} = 4x_{l}\,\delta$ plus a constant field, and with (a), $\sum x_{l}\,\tr S = 0$. (c) Finally $Df' = -2\,\tr G[f']$ by (2), so $Df'$ vanishes outside $B$ and has no monopole and no dipole. If $Df' = \delta_{a} - \delta_{b}$ and $G[f']$ vanished outside such a box, then $a$ and $b$ would lie in it (elsewhere $\tr G[f']$ would be nonzero at a site outside the box) and the dipole $a - b$ would be zero. The commutator of the curvature with $O_{f'}$ is $g\,pq\,G[f']\,X_{f'}$ by the computation of (4), which uses only $[h_{c}(x), W(\pm gf')]$. The count for $f' = \phi\,\delta$ is computed.

\emph{(4).} For a matter element $m$, $[C_{0}(s), m\,W(f')] = \bigl((Df')(s)\,m - g\,[\rho_{s}, m]\bigr)W(f')$. The mass density is built from the commuting numbers of the arrangements, as the source of Proposition \claimref{field-theory.operator-newtonian-coupling} is. From $P_{1}v = v = vP_{2}$, $vP_{1} = P_{2}v = 0$ (\claimref{synthesis.set-and-clear-of-a-bit-generate-the-hyperbolic-plane}): $[\rho_{s}, v] = (\rho^{\mathrm{w}} - \rho^{\mathrm{n}})(s)\,v$ and $[\rho_{s}, w] = -(\rho^{\mathrm{w}} - \rho^{\mathrm{n}})(s)\,w$. So $[C_{0}(s), e_{2}] = -g(\rho^{\mathrm{w}} - \rho^{\mathrm{n}})(s)(v + w)$, which is the first statement, and $[C_{0}(s), v\,W(gf')] = g\bigl(Df' - \rho^{\mathrm{w}} + \rho^{\mathrm{n}}\bigr)(s)\,v\,W(gf') = 0$, likewise for $w\,W(-gf')$; the $C_{i}$ commute with matter elements and with every $W$. $X_{f'}^{2} = -(vw + wv) = -1$ since $v^{2} = w^{2} = 0$ and the Weyl operators commute with the matter and with one another; these are the steps of Theorem \claimref{synthesis.two-dimensional-lattice-maxwell-coulomb-slaved-field-fails-and-string-dressed-detection-commutes}(1), (2) with the Hamiltonian constraint in place of the Gauss constraint. The pair $(h, \pi)$ is canonical and so is not an element of a floor of the tower (\claimref{field-theory.no-canonical-commutation-relations-tower}); it enters only through the Weyl operators and their stated relations. $O_{f'} = (p - qX)P_{1}(p + qX) = p^{2}P_{1} + pq(P_{1}X - XP_{1}) - q^{2}XP_{1}X$ with $P_{1}X = v\,W(gf')$, $XP_{1} = -w\,W(-gf')$, $XP_{1}X = -wv = -P_{2}$. From $[h_{c}(x), W(\pm gf')] = \pm g f'_{c}(x)W(\pm gf')$ follows $[h_{c}(x), O_{f'}] = g\,pq\,f'_{c}(x)\,X_{f'}$, and the curvature is linear in $h$. The counts are computed: $6\cdot 729 = 4374$ components, $9$ on the string, $57$ in the support of $G[f]$.

\emph{(5).} For $h = \phi\,\delta$: $G[h]_{ii} = \tfrac12(\Delta - \Delta_{i})\phi$, $G[h]_{ij} = -\tfrac12 s_{i}s_{j}\phi$, $\tr G[h] = \Delta\phi$, $Dh = -2\Delta\phi$. So $G[h] - \tfrac12\delta\,\tr G[h]$ has components $-\tfrac12 b_{i}s_{i}\phi$ and $-\tfrac12 s_{i}s_{j}\phi$: it is the displacement $\delta_{\xi}$ with $\xi_{i} = -\tfrac14 s_{i}\phi$. Starting from $h_{0} = \phi\,\delta$, $\pi = 0$, induction on the step gives $\pi(n) = -n\tau\,G[h_{0}]$ and $h(n) = h_{0} - n(n+1)\tau^{2}\delta_{\xi}$, so $G[h(n)] = G[h_{0}]$. The constraints hold: $D(\phi\,\delta) = -2\Delta\phi$ is the density with its uniform number, and $C_{i}[\pi(n)] = -n\tau\,C_{i}[G[h_{0}]] = 0$. The curvature of the slaved field is $g\,G[\phi^{\mathrm{w}}\delta]P_{1} + g\,G[\phi^{\mathrm{n}}\delta]P_{2} = g\kappa P_{1} + g\,G[\phi^{\mathrm{n}}\delta]$, whence the commutators; $[P_{1}, O] = pq\,e_{2}$ as in Theorem \claimref{synthesis.field-slaved-to-an-operator-source-does-not-commute-at-spacelike-separation}(1). The potentials are obtained exactly (the difference $\phi^{\mathrm{w}} - \phi^{\mathrm{n}}$ from the first linear relation among the iterates of $\Delta$ on its source, verified by applying $\Delta$), and the count $3858$ is computed.

\emph{(6).} For any $h$ let $\phi_{h}$ solve $-2\Delta\phi_{h} = Dh$ (possible since $Dh$ sums to zero). Then $h - \phi_{h}\delta$ has $D = 0$ and $G[h] = G[h - \phi_{h}\delta] + G[\phi_{h}\delta]$. By the symmetry of $G$, the trace identity and the symmetry of $\Delta$: $w_{c}\,G[\phi_{h}\delta]_{c}(x) = \langle\epsilon_{c,x}, G[\phi_{h}\delta]\rangle = \sum_{s}\tr G[\epsilon_{c,x}](s)\,\phi_{h}(s) = -\tfrac12\sum_{s}(\Delta u_{c,x})(s)\,\phi_{h}(s) = \tfrac14\sum_{s}u_{c,x}(s)\,(Dh)(s)$. So $G^{L}[h] = G[\phi_{h}\delta]$, a function of $Dh$, equal on the constraint surface to $g$ times the same expression in $\rho$. Its commutator with $v\,W(gf)$ is $g\,G^{L}[f]$ times it, and $G^{L}[f] = G[\phi_{f}\delta] = \kappa$ since $Df = \rho^{\mathrm{w}} - \rho^{\mathrm{n}}$; the rest follows from (4) by subtraction.

\emph{(7).} $D(g f P_{1} + gF^{\mathrm{n}}) = g(\rho^{\mathrm{w}} - \rho^{\mathrm{n}})P_{1} + g\rho^{\mathrm{n}} = g\rho$ up to the uniform number, for a number field $F^{\mathrm{n}}$ with $DF^{\mathrm{n}} = \rho^{\mathrm{n}} - \bar\rho$ (for instance $\phi^{\mathrm{n}}\delta$), and $\pi = 0$ satisfies $C_{i} = 0$. Its commutators are those of $P_{1}$. That $D$ of the evolved field is unchanged is (2), so its constrained part is $g\kappa P_{1}$ plus a number at every step; the count $499$ is computed.

\emph{(8).} The matter is at rest and the field equations have no source, so with the field's own variables $h(n)$ and $\pi(n)$ are linear in $h(n_{d})$, $\pi(n_{d})$, and conjugating by $W(gf)$ at step $n_{d}$, which shifts $h(n_{d})$ by $gf$ and leaves $\pi(n_{d})$, shifts $h(n)$, $\pi(n)$ by $g$ times the solution $S(n)$ with data $(f, 0)$ at step $n_{d}$ (forward by the step, backward by its inverse). Hence $[h(n), W(\pm gf)] = \pm gS_{h}(n)W(\pm gf)$ and the stated commutator with $O_{f}$ as in (4). A step changes $\pi$ at a component only by $G$ of $h$ on that component's stencil and $h$ only by $\pi$ at the same site, so $S(n)$ vanishes on every component of depth greater than $|n - n_{d}|$, the depth being counted from the nine components on which $f$ is nonzero (the backward steps reach no further than the forward ones), and $G[S(n)]$ vanishes at every event outside the cone. For the slaved field, $h^{\mathrm{w}}$ and $h^{\mathrm{n}}$ obey the field equations and at step $0$ differ by $gf$ with equal $\pi$, so their difference is $gR(n)$; $D$ of each is $g$ times its density at every step by (2). The commutator with $O$ is that of $P_{1}$. $G[R(n)]$ vanishes outside the lattice's cone of the string at step $0$ by the same locality. The counts are computed. So is the split against a cone of unit speed: the lattice wave equation of (2) with $\tau = \tfrac12$ has unit speed in sites per unit time, that is half a site a step, while a step reaches one site; with $d$ the Euclidean distance from the base site of a curvature component to the nearest of the $22$ sites where $G[f] \neq 0$, the distance being the least-image one on the periodic lattice and a component being placed at its base site, the failures with $d \le n/2$ number $3490$ and those with $d > n/2$ number $5246$; of the $3490$, $3423$ have $d > |n - n_{d}|/2$, $604$ of them at steps $n \le n_{d}$, and the other $67$ all have $n = n_{d}$ and $d = 0$. The second group are events where a solution of the continuum equation with the same data would vanish; the lattice solution is nonzero there, exactly, because its domain of dependence is the wider cone.

\emph{(9).} Slaved: $G[h^{\mathrm{sl}}]_{c}(x)(n) = N + g\,G[R(n)]_{c}(x)\,P_{1}$ with $N$ a number, so $V = e^{-\lambda IN}(e^{-I\vartheta}P_{1} + P_{2})$, and the change of the probability is clause (2) of Theorem \claimref{synthesis.with-a-coupled-device-the-slaved-field-failure-is-a-first-order-signal-and-a-survey-finds-no-stored-bound} with this $\vartheta$ (verified here at two rational points of the circle). With the field's own variables, $V$ is the exponential of a linear function of $h(n_{d})$, $\pi(n_{d})$, and $[G[h]_{c}(x)(n), W(\pm gf)] = \pm g\,G[S(n)]_{c}(x)\,W(\pm gf)$ by (8), so $V^{-1}W(\pm gf)V = e^{\pm I_{4}\theta}W(\pm gf)$; $V$ commutes with the matter. Insert in the form of $O_{f}$ in (4). Outside the cone $G[S(n)]_{c}(x) = 0$ by (8).

\emph{(10).} In two space dimensions the symmetric maps of the space directions have dimension three. For nonzero $\kappa$ in the plane the maps $a \mapsto (a\cdot\kappa)\xi + (a\cdot\xi)\kappa$, $\xi$ in the plane, have dimension two, and $G_{\kappa}$ vanishes on them: at $\omega = 0$ they are displacements with $h(e_{0}) = 0$, on which $\mathcal{G}_{\kappa}$ vanishes by the stored theorem, and by (1) at $\omega = 0$ the part of $\mathcal{G}_{\kappa}[h](a)$ orthogonal to $e_{0}$ is $-G_{\kappa}[h](a)$. The remaining direction is $h(a) = (a\cdot n)\,n$ with $n$ a unit vector of the plane orthogonal to $\kappa$, $Q(n) = -1$: then $h(\kappa) = 0$, $\tr h = Q(n) = -1$ and $G_{\kappa}[h](a) = \tfrac12\bigl[(a\cdot\kappa)\kappa - Q(\kappa)(a\cdot n)n - Q(\kappa)\,a\bigr]$, which is zero because $a = (a\cdot\kappa)\kappa/Q(\kappa) - (a\cdot n)n$ in the plane. So $G_{\kappa} = 0$ on all symmetric maps of two space directions (also verified on three sample maps in $\Cl(1,2)$). On the lattice of $16$ sites the ranks are computed: $48$ components, displacements of rank $30$, $D$ of rank $15$, leaving the three constant fields. So at nonzero momentum $h$ modulo displacements is determined by $Dh$, and on the constraint surface $Dh = g\rho$.

\emph{Remark.} $S(n)P_{1} + N(n)$ obeys the field equations, linear and without source, and $D$ of it is $g(\rho^{\mathrm{w}} - \rho^{\mathrm{n}})P_{1}$ plus $D N$ at every step; its commutator with $O$ at step $n_{d}$ is $g\,G[S(n)]\,pq\,e_{2}$, zero outside the cone by (8). The count for a detection at step $5$ is computed.
